% T2 · Por qué conmutar y no disipar — lámina vectorial de pizarra
{\setbeamercolor{background canvas}{bg=pizarra}
\begin{frame}[plain]
\begin{tikzpicture}[piz]
\begin{scope}[shift={(current page.south west)}]
  \ltitulo{T2 · Por qué conmutar y no disipar}

  % ================= modo lineal =================
  \draw[panel] (0.40,4.45) rectangle (7.70,8.05);
  \ptitulo{0.62}{7.98}{Modo lineal: el transistor es una resistencia}
  \begin{scope}[shift={(1.35,5.95)}]
    % fuente
    \draw[tiza,thick] (0,0.30) -- (0,0.62);
    \draw[tiza,thick] (-0.28,0.62) -- (0.28,0.62);
    \draw[tiza,thick] (-0.15,0.80) -- (0.15,0.80);
    \draw[tiza,thick] (0,0.80) -- (0,1.20);
    \node[text=tiza,font=\tiny,anchor=east] at (-0.34,0.71) {5 V};
    % motor
    \draw[tiza,thick] (0,1.20) -- (1.35,1.20);
    \draw[tiza,thick] (1.75,1.20) circle (0.30);
    \node[text=tiza,font=\tiny] at (1.75,1.20) {M};
    \draw[tiza,thick] (2.05,1.20) -- (3.30,1.20) -- (3.30,0.95);
    % elemento lineal
    \draw[tiza,thick] (3.02,0.95) rectangle (3.58,0.32);
    \draw[ambar,-{Stealth[length=4pt]},thick] (4.12,0.85) -- (3.66,0.68);
    \node[text=ambar,font=\tiny,anchor=west] at (4.16,0.85) {2.5 V};
    % retorno
    \draw[tiza,thick] (3.30,0.32) -- (3.30,0) -- (0,0) -- (0,0.30);
    \draw[cian,-{Stealth[length=4pt]}] (1.30,0) -- (1.75,0);
    \node[text=cian,font=\tiny,anchor=south] at (1.55,0.04) {0.40 A};
  \end{scope}
  \node[anchor=north west,text=rojo,font=\sffamily\bfseries\large]
    at (0.62,5.85)
    {$P = 2.5\,\mathrm{V}\times0.40\,\mathrm{A} = 1.0\,\mathrm{W}$};
  \ptexto{0.62}{5.32}{6.9cm}{Toda esa tensión sobrante se vuelve calor
    \emph{todo el tiempo}.}

  % ================= modo conmutado =================
  \draw[panel] (7.95,4.45) rectangle (15.60,8.05);
  \ptitulo{8.17}{7.98}{Modo conmutado: solo dos estados}
  \draw[rot] (11.77,5.05) -- (11.77,7.45);
  \node[text=cian,font=\bfseries\small] at (9.90,7.32) {ON};
  \node[text=cian,font=\bfseries\small] at (13.75,7.32) {OFF};
  \node[anchor=north west,text=tiza,align=left,text width=3.3cm,
    font=\sffamily\scriptsize] at (8.30,7.02)
    {$V_{DS}=I_D\,R_{DS(on)}$\\ $=0.40\times0.15=60\,$mV\\[2pt]
     $I_D=0.40\,$A};
  \node[anchor=north west,text=tiza,align=left,text width=3.3cm,
    font=\sffamily\scriptsize] at (12.05,7.02)
    {$V_{DS}=5\,$V\\[2pt] $I_D\approx0$ (solo fuga)\\[2pt]
     la corriente no circula};
  \node[anchor=north west,text=cian,font=\sffamily\bfseries\large]
    at (8.30,5.85) {$P_{ON}=24\,$mW};
  \node[anchor=north west,text=cian,font=\sffamily\bfseries\large]
    at (12.05,5.85) {$P_{OFF}\approx0$};
  \node[anchor=north west,text=ambar,align=left,text width=7.2cm,
    font=\sffamily\scriptsize] at (8.17,5.30)
    {\textbf{40 veces menos calor} que el modo lineal, con el mismo
     motor y la misma corriente.};

  % ================= dónde se paga =================
  \draw[panel] (0.40,0.40) rectangle (15.60,4.20);
  \ptitulo{0.62}{4.08}{Pero el flanco no es gratis}
  % ejes
  \draw[tizasuave,-{Stealth[length=4pt]}] (0.95,1.60) -- (7.30,1.60)
    node[right,font=\tiny] {$t$};
  % v_DS
  \draw[tiza,thick] (1.05,3.35) -- (2.20,3.35) -- (3.00,1.72)
    -- (5.10,1.72) -- (5.90,3.35) -- (7.10,3.35);
  \node[text=tiza,font=\tiny,anchor=east] at (1.02,3.35) {$v_{DS}$};
  % i_D
  \draw[cian,thick] (1.05,1.68) -- (2.20,1.68) -- (3.00,2.90)
    -- (5.10,2.90) -- (5.90,1.68) -- (7.10,1.68);
  \node[text=cian,font=\tiny,anchor=east] at (1.02,1.90) {$i_D$};
  % solape de potencia
  \fill[ambar,opacity=0.40] (2.20,1.60) -- (2.60,2.35) -- (3.00,1.60) -- cycle;
  \fill[ambar,opacity=0.40] (5.10,1.60) -- (5.50,2.35) -- (5.90,1.60) -- cycle;
  \node[text=ambar,font=\tiny,anchor=east] at (2.15,2.05) {$p=v\,i$};
  \node[text=ambar,font=\tiny,anchor=south] at (4.05,2.10)
    {el costo se paga en los flancos};
  \node[anchor=north west,text=tizasuave,align=left,text width=6.3cm,
    font=\sffamily\tiny] at (0.95,1.42)
    {Modelo simplificado con carga resistiva. Con carga inductiva el
     solape es mayor y depende del camino de recirculación.};
  % términos de pérdida
  \node[anchor=north west,text=tiza,align=left,text width=7.5cm,
    font=\sffamily\footnotesize] at (7.85,3.72)
    {$P_{cond}=I_D^{2}R_{DS(on)}=24\,$mW \quad
     \textcolor{tizasuave}{no depende de $f$}\\[5pt]
     $P_{sw}=(E_{on}+E_{off})\,f$ \quad
     \textcolor{ambar}{crece en proporción a $f$}\\[5pt]
     $P_{fuga}=V_{DS}\,I_{DSS}$ \quad
     \textcolor{tizasuave}{despreciable aquí, no siempre}};
  \node[anchor=north west,text=cian,align=left,text width=7.5cm,
    font=\sffamily\footnotesize] at (7.85,2.05)
    {Subir $f$ reduce rizado y filtro, pero se paga en cada flanco. Ahí
     $Q_g$ deja de ser un detalle del datasheet y se vuelve el límite
     del diseño.};
\end{scope}
\end{tikzpicture}
\end{frame}}
